\section{The real Bring radical as a primitive-recursive real}\label{sec:real-bring}

Set
\[
 P_a(x)=x^5+x+a,
 \qquad h(x)=x^5+x.
\]
The usual proof that \(h\) is strictly increasing appeals to ordinary real
calculus.  For \ipra{} it is better to use a polynomial certificate that is
valid over every ordered field.

\subsection{A sum-of-squares monotonicity certificate}

Define
\[
 S_4(x,y)=x^4+x^3y+x^2y^2+xy^3+y^4.
\]

\begin{lemma}[Positivity of the fifth-power divided difference]
\label{lem:s4-sos}
The following identity holds in \(\Z[x,y]\):
\begin{equation}\label{eq:s4-sos}
 12S_4(x,y)
 =3(2x^2+xy)^2+(3xy+2y^2)^2+8y^4.
\end{equation}
Consequently \(S_4(x,y)\ge0\) in every ordered field.
\end{lemma}

\begin{proof}
Expansion of the right side gives
\[
 12x^4+12x^3y+12x^2y^2+12xy^3+12y^4.
\]
This is an exact polynomial normalization, so after coding signed coefficients
it is an \ipra{} ring identity.  Nonnegativity follows from nonnegativity of
squares and multiplication by the positive integer \(12\).
\end{proof}

\begin{theorem}[Strong monotonicity]\label{thm:strong-monotonicity}
For \(y\ge x\),
\begin{equation}\label{eq:strong-monotonicity}
 P_a(y)-P_a(x)
 =(y-x)(1+S_4(x,y))\ge y-x.
\end{equation}
In particular \(P_a\) and \(h\) are strictly increasing.
\end{theorem}

\begin{proof}
The factorization
\(y^5-x^5=(y-x)S_4(x,y)\) is exact.  Apply
Lemma~\ref{lem:s4-sos}.  If \(y>x\), then the right side is at least the
positive number \(y-x\).
\end{proof}

The extra lower bound in~\eqref{eq:strong-monotonicity}, rather than mere strict
increase, is the key constructive estimate used throughout this section.

\subsection{Exact rational bisection}

Let \(a\in\Q\).  Put
\[
 M=\lceil |a|\rceil+1.
\]
Then
\[
 P_a(-M)<0<P_a(M).
\]
Starting with \([L_0,U_0]=[-M,M]\), define
\(m_n=(L_n+U_n)/2\) and
\begin{equation}\label{eq:bisection}
 (L_{n+1},U_{n+1})=
 \begin{cases}
   (m_n,U_n),&P_a(m_n)\le0,\\
   (L_n,m_n),&P_a(m_n)>0.
 \end{cases}
\end{equation}
Every comparison in~\eqref{eq:bisection} is a decidable comparison of rational
codes.

\begin{lemma}[Bisection invariant]\label{lem:bisection-invariant}
For every \(n\),
\[
 L_n\le U_n,
 \qquad P_a(L_n)\le0\le P_a(U_n),
 \qquad U_n-L_n=2M\,2^{-n}.
\]
The functions \(L_n,U_n\) are primitive recursive in the code of \(a\) and
\(n\).
\end{lemma}

\begin{proof}
One step of~\eqref{eq:bisection} preserves the two endpoint signs by the branch
condition and halves the width.  All statements are rational equalities or
orders and the induction invariant can be Boolean-folded, so quantifier-free
induction suffices.
\end{proof}

Let
\[
 N(a,k)=k+2+\left\lceil\log_2(2M)\right\rceil,
 \qquad
 q_a(k)=\frac{L_{N(a,k)}+U_{N(a,k)}}2.
\]
Then the containing interval has width at most \(2^{-k-1}\).  The sequence
\(q_a\) is therefore a regular dyadic name.

\begin{theorem}[Constructive existence and uniqueness]\label{thm:bring-exists}
For every rational \(a\), the name \(q_a\) converges with its displayed
primitive-recursive modulus to a real \(x\) satisfying
\[
 x^5+x+a=0.
\]
This zero is unique.
\end{theorem}

\begin{proof}
Nested intervals with a geometric width give a regular Cauchy name.  On
\([-M,M]\), the exact factorization of fifth powers gives the rational
Lipschitz estimate
\[
 |P_a(u)-P_a(v)|
 \le \bigl(1+5M^4\bigr)|u-v|.
\]
Since \(P_a(L_n)\le0\le P_a(U_n)\), monotonicity places
\(P_a(q_a(k))\) between the endpoint values; the difference of those endpoint
values tends effectively to zero by the displayed Lipschitz bound.  Hence the
limit satisfies \(P_a(x)=0\).  If \(x<y\) were two zeros,
Theorem~\ref{thm:strong-monotonicity} would give
\(0=P_a(y)-P_a(x)\ge y-x>0\), a contradiction.
\end{proof}

\begin{definition}[Real Bring radical]
For rational \(a\), \(\BR(a)\) is the regular name constructed above.  For a
regular real name \(a_n\), define \(\BR(a)\) by applying rational bisection to
a sufficiently accurate \(a_n\) and using the stability estimate in the next
theorem.  This extension is independent of the chosen name.
\end{definition}

\subsection{Stability, oddness, and order}

\begin{theorem}[Sharp inverse Lipschitz estimate]\label{thm:bring-lipschitz}
For all real \(a,b\),
\begin{equation}\label{eq:bring-lipschitz}
 \abs{\BR(a)-\BR(b)}\le |a-b|.
\end{equation}
More generally, for every real \(q\),
\begin{equation}\label{eq:residual-bound}
 \abs{q-\BR(a)}\le |q^5+q+a|.
\end{equation}
\end{theorem}

\begin{proof}
Let \(x=\BR(a)\), \(y=\BR(b)\), and suppose \(x\le y\).  Strong monotonicity
gives
\[
 y-x\le h(y)-h(x)=(-b)-(-a)=a-b\le|a-b|.
\]
The other order is symmetric.  Taking \(y=q\) and using
\(P_a(x)=0\) gives~\eqref{eq:residual-bound}.
\end{proof}

Estimate~\eqref{eq:bring-lipschitz} proves at once that applying rational
bisection to successively better rational approximants of a named real gives a
regular name.  This is the constructive extension promised above.

\begin{corollary}[Oddness and strict decrease]\label{cor:bring-odd-decreasing}
The Bring radical is odd and strictly decreasing:
\[
 \BR(-a)=-\BR(a),
 \qquad a<b\Longrightarrow \BR(b)<\BR(a).
\]
\end{corollary}

\begin{proof}
If \(x^5+x+a=0\), then \((-x)^5+(-x)-a=0\); uniqueness gives oddness.  If
\(a<b\), then
\[
 P_b(\BR(a))=b-a>0.
\]
Since \(P_b\) is strictly increasing and vanishes at \(\BR(b)\), its positive
value at \(\BR(a)\) forces \(\BR(b)<\BR(a)\).
\end{proof}

\begin{corollary}[Unboundedness with a witness]\label{cor:bring-unbounded}
For every \(m\ge0\),
\[
 a>m^5+m\Longrightarrow \BR(a)<-m,
 \qquad
 a<-(m^5+m)\Longrightarrow \BR(a)>m.
\]
Thus \(\BR\) is unbounded in both directions.
\end{corollary}

\begin{proof}
For the first implication,
\(P_a(-m)=a-m^5-m>0\).  The zero of the increasing polynomial lies to the
left of \(-m\).  The second implication follows by oddness.
\end{proof}

\subsection{Differentiability without an inverse-function black box}

Let \(x=\BR(a)\), \(y=\BR(a+t)\).  Subtracting the two defining equations gives
the exact identity
\begin{equation}\label{eq:bring-difference-quotient}
 (y-x)\bigl(1+S_4(x,y)\bigr)=-t.
\end{equation}
The denominator is at least one by Lemma~\ref{lem:s4-sos}.

\begin{theorem}[Derivative of the real Bring radical]\label{thm:bring-derivative}
The real Bring radical is effectively differentiable and
\begin{equation}\label{eq:bring-derivative}
 \BR'(a)=-\frac{1}{1+5\BR(a)^4}.
\end{equation}
On every rationally bounded interval, the convergence of the difference
quotient has a primitive-recursive modulus.
\end{theorem}

\begin{proof}
By~\eqref{eq:bring-lipschitz}, \(|y-x|\le|t|\).  Hence
\(S_4(x,y)\to S_4(x,x)=5x^4\) with an explicit polynomial modulus on any
supplied bound for \(|x|\).  Divide~\eqref{eq:bring-difference-quotient} by
\(t\ne0\) and pass to the effective limit.  Since the denominator is uniformly
at least one, reciprocal error propagation is primitive recursive.
\end{proof}

Repeated differentiation produces rational functions of \(\BR(a)\) whose
denominators are powers of \(1+5\BR(a)^4\).  For example,
\[
 \BR''(a)=-\frac{20\BR(a)^3}{(1+5\BR(a)^4)^3}.
\]
Thus each fixed derivative is a named continuous function.  A quantitative
complex inverse theorem applied at a real point gives a local analytic branch;
all constants are computable because \(5x^4+1\ge1\) on the real axis.  This
establishes, in the effective sense, the source page's assertion that the
complex branch can be chosen analytic in a neighborhood of the real line.

\subsection{Fifth-root asymptotics}

For \(a>0\), put \(y=-\BR(a)>0\).  Then
\begin{equation}\label{eq:y-equation}
 y^5+y=a.
\end{equation}
Let \(r=a^{1/5}\), represented by its standard primitive-recursive bisection
name.  Equation~\eqref{eq:y-equation} gives \(y\le r\).  Factoring the
difference of fifth powers yields
\begin{equation}\label{eq:asymptotic-exact}
 (r-y)(r^4+r^3y+r^2y^2+ry^3+y^4)=y.
\end{equation}

\begin{theorem}[Effective asymptotic bound]\label{thm:bring-asymptotic}
For \(a=r^5\) with \(r\ge1\),
\begin{equation}\label{eq:bring-asymptotic-bound}
 0\le r+\BR(a)\le r^{-3},
 \qquad
 0\le 1-\frac{-\BR(a)}{r}\le r^{-4}.
\end{equation}
Consequently \(\BR(a)\sim-a^{1/5}\) as \(a\to+\infty\), with the explicit
modulus supplied by~\eqref{eq:bring-asymptotic-bound}.  The negative direction
follows by oddness.
\end{theorem}

\begin{proof}
The parenthesized factor in~\eqref{eq:asymptotic-exact} is at least \(r^4\),
and \(y\le r\).  Therefore
\(0\le r-y\le y/r^4\le r^{-3}\).  Divide by \(r>0\) for the relative bound.
All inequalities become integer polynomial inequalities after multiplying by
positive powers of \(r\), so the proof is compatible with the regular-real
coding.
\end{proof}
